\chapter{Appendix}

This appendix provides code and data availability, feature and method reference
material, implementation and reproducibility details, and supplementary
evaluation output. Its sections cover code availability
(Appendix~\ref{appendix:code-availability}), health indicators
(Appendix~\ref{appendix:health-indicators}), implementation detail
(Appendix~\ref{appendix:impl-detail}), evaluation instruments
(Appendix~\ref{appendix:impl-eval-instruments}), and undefined-feature handling
(Section~\ref{sec:appendix-undefined-features}) for the implementation described
in Chapter~\ref{chap:implementation}.

\section{Code and Data Availability}
\label{appendix:code-availability}

The artifact of this thesis --- the historical dataset builder, the feature and label
computation, the model implementations, the training notebook, the scoring service, and
the orchestration and analysis scripts --- is publicly available at

\begin{center}
\url{https://github.com/esuarkeN/oss-risk-radar}
\end{center}

\noindent The results reported here were produced from the pinned revision
\texttt{f5e688e} \cite{ossRiskRadarReadme}, which carries the reported dataset, the
staged model artifacts, and the training notebook in the state they were evaluated in.
That revision is what the dataset hash in Appendix~\ref{sec:method-reproducibility}
identifies, and all file paths given in this thesis are relative to the repository root.
The commands that regenerate the foundation seed, rebuild the dataset, and retrain the
models are listed in Appendix~\ref{sec:method-reproducibility}, together with the
per-run documentation items. Analyses added after the pinned revision --- among them the
label-definition sensitivity harness of Appendix~\ref{appendix:label-sensitivity} ---
are available on the default branch rather than in \texttt{f5e688e}.

\section{Observable Health Indicator Categories}
\label{appendix:health-indicators}

Table~\ref{tab:background-health-indicators} lists the observable health
categories, example signals, public sources, and expected relevance
(Section~\ref{sec:background-oss-health}). Section~\ref{sec:method-feature-engineering}
implements the corresponding feature groups, and
Appendix~\ref{appendix:feature-reference} defines each feature.
Security-practice indicators remain parallel triage context, not inactivity-model
inputs (Section~\ref{sec:related-security-practice}).

\begin{table}[htbp]
\centering
\caption{Categories of observable health indicators and their expected relevance}
\label{tab:background-health-indicators}
\small
\begin{tabular}{p{0.18\textwidth} p{0.29\textwidth} p{0.19\textwidth} p{0.21\textwidth}}
\hline
\textbf{Category} & \textbf{Example signals} & \textbf{Public source} & \textbf{Expected relevance} \\
\hline
Commit activity & Commits over recent windows; active commit months; days since last commit & GH Archive push events & Direct evidence of ongoing development \\
Contributor activity & Recent and new contributors; contributor concentration & GH Archive commit actors & Team breadth and single-maintainer (bus-factor) risk \\
Issue and PR activity & Opened/closed issues; PR merge ratio; response and merge latency & GH Archive issue/PR events & Maintenance interaction and responsiveness \\
Release activity & Releases; days since last release; versions published & GH Archive releases and package registry & Whether usable versions still reach users \\
Age and popularity & Repository and package age; stars; forks; popularity tier & GitHub metadata and package registry & Maturity and reach context \\
Security practice & Selected OpenSSF Scorecard checks & OpenSSF Scorecard & Parallel security-posture context, not an inactivity predictor \\
\hline
\end{tabular}
\end{table}

\section{Full-History Feature Reference}
\label{appendix:feature-reference}

This section expands the feature engineering step of Section~\ref{sec:method-feature-engineering}.

The full-history model consumes a feature vector of 60 variables under feature set \texttt{feature\sdash set\sdash v4\sdash full\sdash history}. These come from two sources, documented separately below: 43 \emph{reconstructed historical} features, derived from the GH Archive event history before \(t\) and summarized by group in Table~\ref{tab:method-feature-groups} and individually in Table~\ref{tab:method-feature-reference}; and 17 \emph{current-snapshot} features, obtainable from a single observation of the repository and package metadata, listed in Table~\ref{tab:method-snapshot-features}. The cold-start regime uses exactly the 17 current-snapshot features and none of the historical ones, which is what makes it applicable to repositories with no reconstructed history. Table~\ref{tab:method-feature-reference} additionally lists two variables that are computed and exported for inspection but are \emph{not} part of the 60-variable model vector: \texttt{repo\_archived\_at\_obs}, excluded so that archival cannot become a shortcut label, and \texttt{repo\_age\_days}, which is identically zero under the reproducible offline build because that build carries no repository creation date (Section~\ref{sec:method-offline-metadata}). Because the feature names are compact identifiers whose exact meaning is not obvious to a reader who is new to the artifact, this appendix documents each feature individually. The identifiers used here are the same strings that appear in the training code, in the exported model artifacts, and in the runtime interface, so the tables double as a shared glossary across the thesis and the software.

Four conventions apply throughout and explain the recurring suffixes in the names. First, a numeric suffix denotes a trailing time window that ends at the observation time \(t\): \texttt{\_30d}, \texttt{\_90d}, and \texttt{\_365d} count events in the 30, 90, and 365 days immediately before \(t\), and \texttt{\_at\_obs} denotes a state measured exactly at \(t\). Second, all commit- and contributor-based counts use a human-actor filter, so activity generated by bots does not inflate the maintenance evidence. Third, wherever an absolute count would not be comparable across projects of very different sizes, the feature set also provides a relative or trend-based form (a ratio, a share, or a year-over-year change), because a large project and a small project can be equally well or badly maintained at very different absolute volumes. Fourth, a feature that is \emph{undefined} for a repository is distinguished from one that is genuinely zero (Section~\ref{sec:method-undefined-features}): the ``days since'' features are right-censored at the age of the observable record, and the latency medians are left undefined and imputed by the model artifact rather than being set to zero, which would read as an instantaneous response. Every value is computed from information available at or before \(t\) only; the reasoning behind this restriction, and the distinction between it and the legitimate activity-persistence effect, is discussed in Section~\ref{sec:method-leakage-prevention}.

\begin{table}[htbp]
\centering
\caption{Implemented full-history feature groups}
\label{tab:method-feature-groups}
\begin{tabular}{p{0.25\textwidth} p{0.65\textwidth}}
\hline
\textbf{Feature group} & \textbf{Example features} \\
\hline
Dependency context &
Dependency count of the resolved package version at the observation time. \\
Commit activity &
Commits in the last 30, 90, and 365 days; active commit months; days since last commit. \\
Contributor activity &
Contributors in the last 90 and 365 days; new contributors; contributor concentration. \\
Issue activity &
Opened and closed issues; issue closure ratio; issue backlog growth; stale open issues. \\
Pull request activity &
Opened, merged, and closed-unmerged pull requests; pull request merge ratio; stale open pull requests; pull-request response and merge latency. \\
Release activity &
Releases in the last 365 days; days since last release; versions published in the last 365 days; recent release flag. \\
Age and popularity proxies &
Package age; repository age (diagnostic only); stars and forks observed in the available historical data; popularity tier. \\
Derived risk proxies &
Activity drop; contributor drop; release gap risk; concentration risk score. \\
Responsiveness and backlog pressure &
Issue first-response time; issue resolution time; stale issue share; pull-request response and merge latency. \\
\hline
\end{tabular}
\end{table}

{\footnotesize
\setlength{\tabcolsep}{4pt}
\begin{longtable}{p{0.23\textwidth} p{0.40\textwidth} p{0.29\textwidth}}
\caption[Historical feature reference]{Reconstructed historical feature reference
(\texttt{feature-set-v4-full-history}): exact definition and rationale for each
of the 43 historical model inputs, plus the two diagnostic-only variables
\texttt{repo\_archived\_at\_obs} and \texttt{repo\_age\_days}.}
\label{tab:method-feature-reference} \\
\hline
\textbf{Feature} & \textbf{Definition (all measured at or before \(t\))} & \textbf{Why it is included} \\
\hline
\endfirsthead
\multicolumn{3}{l}{\small\itshape Table~\ref{tab:method-feature-reference} continued from previous page} \\
\hline
\textbf{Feature} & \textbf{Definition (all measured at or before \(t\))} & \textbf{Why it is included} \\
\hline
\endhead
\hline
\multicolumn{3}{r}{\small\itshape continued on next page} \\
\endfoot
\hline
\endlastfoot

\multicolumn{3}{l}{\textit{Commit activity}} \\
\hline
\feat{commits\su 30d} & Number of human commits in the 30 days before \(t\). & Short-term development pulse; captures whether the project is moving right now. \\
\feat{commits\su 90d} & Number of human commits in the 90 days before \(t\). & Quarter-scale activity level that is less noisy than the 30-day count. \\
\feat{commits\su 365d} & Number of human commits in the 365 days before \(t\). & Annual development volume; also the baseline against which the year-over-year drop is measured. \\
\feat{active\su commit\su months\su 365d} & Count of distinct calendar months (0--12) in the last year that contain at least one human commit. & Separates steady maintenance from a single burst: 300 commits in one month is a weaker health signal than the same work spread across ten months. \\
\feat{days\su since\su last\su commit} & Days between \(t\) and the most recent human commit; right-censored at the age of the observable record when no commit is ever observed. & Direct staleness signal; a long silence is one of the strongest indicators of impending inactivity. \\
\hline
\multicolumn{3}{l}{\textit{Contributor activity and concentration}} \\
\hline
\feat{contributors\su 90d} & Distinct human contributors in the last 90 days. & Breadth of recent participation; a proxy for the size of the active team. \\
\feat{contributors\su 365d} & Distinct human contributors in the last 365 days. & Annual contributor base; the reference for the contributor drop. \\
\feat{new\su contributors\su 365d} & Contributors active in the last year who never committed before the window started. & Measures inflow of fresh maintainers; a project with no newcomers is more exposed to attrition. \\
\feat{top1\su contributor\su commit\su share\su 365d} & Share (0--1) of the last year's commits made by the single most active contributor. & Key-person dependency: a high share means the project is fragile if that one person leaves. \\
\feat{top2\su contributor\su commit\su share\su 365d} & Combined commit share of the two most active contributors in the last year. & Detects a very small core team even when no single person dominates. \\
\feat{contributor\su concentration\su index} & Herfindahl index: sum of squared per-contributor commit shares over the last year. & A single scalar for how concentrated contributions are (1.0 = one person, near 0 = evenly spread). \\
\feat{maintainer\su concentration\su flag} & 1 if the top contributor share \(\geq 0.7\) or the top-two share \(\geq 0.85\), else 0. & Binary bus-factor alarm for the clearly single-maintainer case. \\
\hline
\multicolumn{3}{l}{\textit{Issue activity and backlog}} \\
\hline
\feat{opened\su issues\su 90d} & Issues created in the last 90 days. & Inbound user demand and engagement with the project. \\
\feat{closed\su issues\su 90d} & Issues closed in the last 90 days. & Maintainer throughput on the issue tracker. \\
\feat{issue\su closure\su ratio\su 90d} & Closed divided by opened issues over the last 90 days. & Whether maintainers keep pace with inflow, expressed relatively so it is comparable across project sizes. \\
\feat{issue\su backlog\su growth\su 90d} & Relative change in the open-issue count between \(t-90\)\,days and \(t\). & A backlog that is trending upward signals that maintainers are losing control; size-normalized. \\
\feat{stale\su open\su issues\su count\su at\su obs} & Issues opened before \(t-90\)\,days that are still open at \(t\). & Long-unaddressed items indicate neglect of the tracker. \\
\hline
\multicolumn{3}{l}{\textit{Pull-request activity}} \\
\hline
\feat{opened\su prs\su 90d} & Pull requests created in the last 90 days. & Contribution inflow from outside the core team. \\
\feat{merged\su prs\su 90d} & Pull requests merged in the last 90 days. & Direct evidence that maintainers are still integrating outside work. \\
\feat{closed\su unmerged\su prs\su 90d} & Pull requests closed without merging in the last 90 days. & A high count relative to merges can signal rejection or disengagement from contributions. \\
\feat{pr\su merge\su ratio\su 90d} & Merged divided by opened pull requests over the last 90 days. & Responsiveness to contributors, expressed relatively. \\
\feat{stale\su open\su prs\su count\su at\su obs} & Pull requests opened before \(t-90\)\,days that are still open at \(t\). & Ignored contributions discourage further contribution and indicate maintainer absence. \\
\hline
\multicolumn{3}{l}{\textit{Release activity}} \\
\hline
\feat{releases\su 365d} & GitHub releases published in the last 365 days. & Shipping cadence: whether usable versions still reach users. \\
\feat{days\su since\su last\su release} & Days since the most recent release or package version; right-censored at the age of the observable record when no release is ever observed. & Release staleness; complements commit staleness for projects that ship rarely. \\
\feat{versions\su published\su 365d} & Package-registry versions published in the last 365 days. & Package-level shipping, which can differ from repository releases. \\
\hline
\multicolumn{3}{l}{\textit{Dependency context}} \\
\hline
\feat{dependency\su count\su at\su obs} & Declared dependency count of the resolved package version at \(t\). & Proxy for maintenance burden and integration complexity. \\
\hline
\multicolumn{3}{l}{\textit{Age and popularity proxies}} \\
\hline
\feat{package\su age\su days} & Days from the first published package version to \(t\). & Maturity context: young and long-established packages have different baseline inactivity risk. \\
\feat{repo\su age\su days} & Days from repository creation to \(t\). & Diagnostic metadata only, and not a model input: the offline build supplies no creation date, so this column is identically zero (Section~\ref{sec:method-offline-metadata}). \\
\feat{stars\su total\su at\su obs} & Cumulative stars observed up to \(t\) (a lower-bound proxy reconstructed from archive events). & Popularity and visibility; more attention was expected to correlate with continued maintenance, although the association observed in this dataset runs the other way (Section~\ref{sec:results-attribution}). \\
\feat{forks\su total\su at\su obs} & Cumulative forks observed up to \(t\). & Reuse and derivative activity around the project. \\
\feat{popularity\su tier\su at\su obs} & Low / medium / high tier (0/1/2) derived from star and fork thresholds. & Coarse popularity bucket, because inactivity dynamics differ between obscure and widely used projects. \\
\hline
\multicolumn{3}{l}{\textit{Diagnostic and binary flags}} \\
\hline
\feat{repo\su archived\su at\su obs} & 1 if the repository is archived by \(t\), else 0. & Diagnostic metadata only, and not a model input; already-archived rows are excluded from fitting so archival cannot become a shortcut label. \\
\feat{has\su recent\su release\su flag} & 1 if any release or package version appeared in the last 365 days. & Simple ``is it still shipping'' indicator that is robust when exact counts are sparse. \\
\feat{has\su recent\su pr\su merge\su flag} & 1 if any pull request was merged in the last 90 days. & Simple ``are maintainers still integrating'' indicator. \\
\feat{has\su issue\su activity\su 365d} & 1 if at least one issue was opened in the last 365 days, else 0. & Makes ``no issue tracker activity at all'' an explicit signal, so its absence cannot reach the model disguised as a zero-valued response latency. \\
\feat{has\su pr\su activity\su 365d} & 1 if at least one pull request was opened in the last 365 days, else 0. & The pull-request counterpart; distinguishes a project nobody contributes to from one that merges contributions instantly. \\
\hline
\multicolumn{3}{l}{\textit{Derived risk proxies}} \\
\hline
\feat{activity\su drop\su 365d\su vs\su prev\su 365d} & Relative decline in commits from the previous year to the last year (positive = fewer commits). & Captures deceleration, enabling earlier detection of decline than the absolute level alone. \\
\feat{contributors\su drop\su 365d\su vs\su prev\su 365d} & Relative decline in the contributor count from the previous year to the last year. & Detects a shrinking maintainer base before commits fully stop. \\
\feat{release\su gap\su risk} & Score in [0,1] for how overdue the next release is relative to the project's own release cadence (age-based when no release history exists). & Normalizes release silence against the project's normal rhythm rather than an absolute threshold. \\
\feat{concentration\su risk\su score} & Composite score in [0,1] combining top-1 and top-2 shares, the Herfindahl index, and the concentration flag. & A single interpretable bus-factor risk value that consolidates the concentration features. \\
\hline
\multicolumn{3}{l}{\textit{Responsiveness and backlog pressure}} \\
\hline
\feat{issue\su first\su response\su median\su days\su 365d} & Median days from issue creation to the first maintainer response over the last year; \emph{undefined} if no issue received a response in the window. & Slowing first responses indicate disengaging maintainers, independently of resolution speed. \\
\feat{issue\su resolution\su median\su days\su 365d} & Median days from issue creation to closure over the last year; \emph{undefined} if no issue was closed in the window. & How quickly reported problems are actually resolved. \\
\feat{stale\su issue\su share\su at\su obs} & Stale open issues divided by the current open backlog. & The fraction of the backlog that is old, expressed relatively so it is comparable across sizes. \\
\feat{pr\su response\su median\su days\su 365d} & Median days to the first response, merge, or close on pull requests over the last year; \emph{undefined} if no pull request was responded to. & Responsiveness to contributors, the pull-request counterpart to issue first response. \\
\feat{pr\su merge\su latency\su median\su days\su 365d} & Median days from creation to merge for merged pull requests over the last year; \emph{undefined} if no pull request was merged. & Integration speed for work that is ultimately accepted. \\
\hline
\end{longtable}
}

The remaining 17 model inputs are \emph{current-snapshot} features (Table~\ref{tab:method-snapshot-features}). They require only a single observation of the repository and its package metadata, with no event history, which is precisely why they constitute the entire cold-start feature set (\texttt{feature\sdash set\sdash v4\sdash cold\sdash start}). They are also part of the full-history vector, so a full-history model sees both tables.

\begin{table}[htbp]
\centering
\caption{Current-snapshot feature reference: the 17 features obtainable without event history. These form the complete cold-start vector and are also included in the full-history vector.}
\label{tab:method-snapshot-features}
{\footnotesize
\setlength{\tabcolsep}{4pt}
\begin{tabular}{p{0.24\textwidth} p{0.40\textwidth} p{0.28\textwidth}}
\hline
\textbf{Feature} & \textbf{Definition (measured at \(t\))} & \textbf{Why it is included} \\
\hline
\feat{has\su repository\su mapping} & 1 if the package resolves to a source repository. & Without a repository, all activity signals are unavailable; the model must know this. \\
\feat{is\su direct\su dependency} & 1 if the package is a direct rather than transitive dependency. & Triage context: direct dependencies are more readily replaced. \\
\feat{last\su push\su age\su days} & Days since the most recent push to the default branch; right-censored when no push is observed. & The cold-start counterpart to commit staleness. \\
\feat{last\su release\su age\su days} & Days since the most recent release; right-censored when none is observed. & Shipping staleness without needing release history. \\
\feat{release\su cadence\su days} & Mean days between consecutive releases; \emph{undefined} with fewer than two releases. & The project's normal shipping rhythm, against which silence is judged. \\
\feat{recent\su contributors\su 90d} & Distinct contributors in the last 90 days. & Breadth of the currently active team. \\
\feat{contributor\su concentration} & Commit share of the most active contributor. & Key-person dependency without full history. \\
\feat{open\su issue\su growth\su 90d} & Relative change in the open-issue count over 90 days. & Whether the backlog is trending upward. \\
\feat{pr\su response\su median\su days} & Median days to first response on pull requests; \emph{undefined} if none were responded to. & Responsiveness to contributors. \\
\feat{stars\su log1p} & \(\log(1 + \text{stars})\) at \(t\). & Popularity on a compressed scale; raw counts span several orders of magnitude. \\
\feat{forks\su log1p} & \(\log(1 + \text{forks})\) at \(t\). & Reuse and derivative activity, likewise compressed. \\
\feat{open\su issues\su log1p} & \(\log(1 + \text{open issues})\) at \(t\). & Backlog size, compressed and comparable across project sizes. \\
\feat{ecosystem\su npm}, \feat{ecosystem\su pypi}, \feat{ecosystem\su go}, \feat{ecosystem\su maven}, \feat{ecosystem\su other} & Five mutually exclusive one-hot indicators of the package ecosystem. & Baseline inactivity rates and release conventions differ by ecosystem. \\
\hline
\end{tabular}
}
\end{table}

\section{Foundation Seed Composition}
\label{sec:appendix-seed-composition}

The foundation seed is generated from the GitHub search API over public, non-fork repositories with a license object and at least 100 stars, drawn separately from three sampling buckets: active, dormant, and archived repositories. The buckets are sampling provenance only and never enter the supervised label (Section~\ref{sec:method-data-sources}). Table~\ref{tab:appendix-seed-strata} contrasts the target frame with the seed as actually built. Two buckets returned fewer repositories than requested --- the GitHub search frame simply did not contain enough matching dormant and low-star active projects --- and a fallback bucket made up the missing count with active repositories instead. That is why the realized seed is majority-active despite a frame designed to oversample inactivity.

\begin{table}[htbp]
\centering
\caption{Foundation seed buckets: target frame versus realized seed (5{,}000 repositories).}
\label{tab:appendix-seed-strata}
\begin{tabular}{lrrrr}
\toprule
Sampling bucket & Target & Realized & Target share & Realized share \\
\midrule
Active   & 2{,}250 & 2{,}601 & 45.0\% & 52.0\% \\
Dormant  & 1{,}750 & 1{,}399 & 35.0\% & 28.0\% \\
Archived & 1{,}000 & 1{,}000 & 20.0\% & 20.0\% \\
\midrule
Total    & 5{,}000 & 5{,}000 & 100\% & 100\% \\
\bottomrule
\end{tabular}
\end{table}

Each of the three sampling groups is drawn from several star-range buckets, which is why the shortfalls are visible only at bucket level. The realized active count of 2{,}601 comprises 2{,}095 repositories from four active buckets --- of which the \texttt{active\sdash foundation} bucket (100--499 stars) returned 845 of the 1{,}000 requested --- plus a 506-repository fallback bucket used to reach the target size. The dormant count of 1{,}399 comprises three buckets, the \texttt{dormant\sdash foundation} bucket returning 649 of 1{,}000; the three archived buckets returned their full 1{,}000.\kiref{Claude Opus 5/8}

\section{Undefined Feature Handling}
\label{sec:appendix-undefined-features}

This section gives the exact rules summarized in Section~\ref{sec:method-undefined-features}.

\paragraph{Censoring bound.} The age of the observable record is the interval from the later of the repository creation date and the start of archive coverage up to the observation time \(t\). When a repository has no observed commit (or no observed release), the corresponding ``days since'' feature is set to this bound, which is the largest value consistent with the data. When the record begins at or after \(t\) --- that is, when nothing at all was observable before the observation date --- there is no interval over which the event could have been missed and therefore no bound to censor against; the feature is then undefined rather than zero. In the reported dataset this applies to \(7{,}687\) of \(50{,}000\) rows, whose repositories have no GH Archive event before their observation date.

\paragraph{Undefined values.} The four responsiveness medians (\texttt{issue\_first\_response\_median\_days\_365d}, \texttt{issue\_resolution\_median\_days\_365d}, \texttt{pr\_response\_median\_days\_365d}, \texttt{pr\_merge\_latency\_median\_days\_365d}) and the two censored age features are exported as absent keys rather than as \texttt{null} or \texttt{0}, so that no consumer can silently read them as zero. Downstream, the gradient-boosted trees route undefined values natively through XGBoost's missing-value branch, while the logistic-regression and neural-network models substitute the per-feature training-set median, which is stored in the model artifact alongside the standardization parameters and re-applied identically at inference. Because a substituted median standardizes to approximately zero, an undefined feature contributes little to the linear score rather than pulling it toward the favorable extreme.

\paragraph{Why it matters.} Before this correction, an undefined latency was stored as \(0.0\). Since \(0\) days is the best possible response time, this encoded the absence of any observed response as an instantaneous one. The affected share is large: each of the four medians is undefined for between \(46.8\%\) and \(60.4\%\) of the \(50{,}000\) snapshots, and \(36.7\%\) of snapshots have neither issue nor pull-request activity in the trailing year at all.\kiref{Claude Opus 5/9} The variable was consequently a strong predictor of inactivity with an inverted meaning, and the feature attributions of Section~\ref{sec:results-attribution} reflected data availability rather than maintenance behavior. The same defect made \texttt{days\_since\_last\_commit} non-monotonic, which in turn depressed the trivial recency baseline against which the learned models are compared (Section~\ref{sec:results-baselines}).

\paragraph{Offline-build consequences.} Because the reported dataset is built through the reproducible offline path (Section~\ref{sec:method-offline-metadata}), the repository creation date, archival timestamp, and fork flag are never populated, and three consequences follow. First, \texttt{repo\_age\_days} is identically zero and is therefore excluded from the model vector and from any age analysis. Second, \texttt{repo\_archived\_at\_obs} is identically zero, so the archival condition of the maintenance rule never fires (Section~\ref{sec:method-definitions}) and the rule excluding already-archived rows from fitting removes none. Third, build-time fork exclusion is inert, though the seed is generated with a non-fork filter so forks are already absent. Populating these attributes would require live enrichment of every seed repository and is noted as future work in Chapter~\ref{chap:conclusion}.

\paragraph{Zero-variance inputs in the reported dataset.} A related consequence deserves to be stated explicitly, because it qualifies the feature counts used throughout this thesis. The foundation seed is repository-first: it is drawn from GitHub search rather than from a package registry, and the offline build performs no deps.dev lookup. Several package-derived inputs of the feature vector therefore have nothing to vary over, and take the same value in all \(50{,}000\) rows. Eight of the \(60\) model inputs are constant in this way:

\begin{itemize}
  \item the five ecosystem indicators, since every row carries the same repository-first provenance (\texttt{ecosystem\_other} is \(1\), the other four \(0\));
  \item \texttt{is\_direct\_dependency} and \texttt{has\_repository\_mapping}, both identically \(1\), because a seed repository is by construction a resolved, directly named project;
  \item \texttt{dependency\_count\_at\_obs}, identically \(0\).
\end{itemize}

Two further inputs vary across the dataset but carry no \emph{between-repository} information, and are counted separately for that reason. Both follow from one mechanism: when no registry version history is available and the build runs offline, the builder records a single synthetic publication event for every repository, dated to the start of the observation period, \texttt{2023-01-01}. \texttt{package\_age\_days} is then the distance from that date to the observation date and takes exactly ten distinct values --- \(0, 90, 181, \dots, 821\) --- one per quarterly cohort of \(5{,}000\) rows. \texttt{versions\_published\_365d} counts publications in the trailing year and is therefore \(1\) for the four \(2023\) cohorts and \(0\) for the six later ones. Neither is a property of the repository; each is a function of the observation date alone. Ten of the sixty model inputs consequently carry no repository-level signal: the eight constants above and these two. The effective feature vector of the reported run is closer to \(50\) informative inputs than to \(60\). Seven of the eight constants --- the five ecosystem indicators, \texttt{is\_direct\_dependency}, and \texttt{has\_repository\_mapping} --- belong to the current-snapshot set, and neither cohort-dependent input does, so the cold-start vector is correspondingly closer to \(10\) than to \(17\).

A third input is affected indirectly. \texttt{has\_recent\_release\_flag} is set when either an archive release event or a package publication falls in the trailing year, so the synthetic publication forces it to \(1\) for every row of the four \(2023\) cohorts. From \texttt{2024-01-01} onward it is driven by archive release events alone and does vary between repositories, taking value \(0\) on \(21{,}169\) rows overall. It is informative in the later cohorts and uninformative in the earlier ones, and is therefore not counted among the ten.

One consequence matters for the time-ordered split. Because \texttt{package\_age\_days} increases monotonically with the cohort, its test-set value of \(821\) lies outside the range seen during training, whose maximum is \(639\). A tree routes such a value to the branch of its largest training split, and a linear model extrapolates a trend that carries no repository-level meaning.\kiref{Claude Opus 5/11}

The reported ranking metrics are unaffected. A feature with no variance cannot split a tree, so the gradient-boosted models assign it zero gain and ignore it, and a constant column standardizes to zero for the linear and neural models; all eight accordingly carry a fitted coefficient of exactly zero in the staged logistic-regression artifact, so they are inert rather than misleading. The two cohort-dependent inputs are constant \emph{within} the test cohort, since every test row shares the observation date \texttt{2025-04-01}, so they cannot reorder test rows and cannot change ROC-AUC. They could in principle shift all test probabilities by a common amount, which would show in Brier score and expected calibration error; this is not separately quantified. Two things do follow. First, the counts ``60'' and ``17'' should be read as the size of the implemented vector, not as the number of signals the reported evaluation actually had available. Second, it explains an otherwise puzzling absence in the feature attribution of Section~\ref{sec:results-attribution}: no package- or ecosystem-derived variable appears among the influential features, not because ecosystem is uninformative in general, but because this dataset contains no variation in it to learn from. Testing whether ecosystem carries signal would require a package-first sampling frame and is left to future work (Chapter~\ref{chap:conclusion}).

\section{System Components and Technology Selection}
\label{appendix:impl-detail}

This section gives the component inventory and language choices. The following
sections document the build and training pipeline
(Appendix~\ref{appendix:impl-pipeline}), runtime services
(Appendix~\ref{appendix:impl-runtime}), and evaluation instruments
(Appendix~\ref{appendix:impl-eval-instruments}). Table~\ref{tab:impl-components}
summarizes the components from Section~\ref{sec:impl-architecture}.

\begin{table}[htbp]
\centering
\caption{Implemented system components}
\label{tab:impl-components}
\begin{tabular}{p{0.26\textwidth} p{0.16\textwidth} p{0.46\textwidth}}
\hline
\textbf{Component} & \textbf{Technology} & \textbf{Responsibility} \\
\hline
Scoring and training workspace & Python & Historical dataset builder, feature engineering, model training, model artifacts, and the runtime scoring service. \\
Backend API & Go & Analysis orchestration, provider enrichment, durable job processing, and model-artifact scoring calls. \\
Analyst frontend & TypeScript / Next.js & Submission, repository overview, risk views, and machine-learning evaluation dashboards. \\
Shared schemas & TypeScript & Reusable data contracts shared by the frontend. \\
Orchestration scripts & Node.js & Seed generation, dataset building, notebook execution, and artifact staging commands. \\
Deployment assets & Docker / Kubernetes & Container images, Compose, and the Argo CD-managed Kubernetes deployment. \\
\hline
\end{tabular}
\end{table}

\label{sec:impl-technology}
The components use different programming languages, each selected for the part of the system where it is most appropriate.

Python is used for all data engineering and machine learning, including the historical dataset builder, the feature and label computation, the model implementations, and the scoring service. Python was chosen because the libraries this thesis depends on --- \texttt{numpy} for the from-scratch model implementations and \texttt{xgboost} for the gradient-boosted trees --- are native to its ecosystem, and because it remains among the most widely used languages for data work \cite{stackoverflow_python_2025}, which eases reproduction of the analytical core. The scoring service is implemented with FastAPI.

Go is used for the backend orchestration service. The backend is primarily a concurrent input/output workload: it calls several external providers and processes durable jobs. Go's static typing, its straightforward concurrency model, and its suitability for long-running network services make it a good fit for this role.

TypeScript with the Next.js framework is used for the analyst frontend, and TypeScript is also used for the shared data contracts. TypeScript provides type safety for the user interface and shares type definitions with the backend contracts. The frontend additionally relies on the broader JavaScript and Node.js tooling ecosystem, which is among the most widely used on GitHub \cite{github_language_stats}. This is a tooling choice rather than a statement about the ecosystems the artifact scores: the dataset seed and the demonstrator in Chapter~\ref{chap:demonstrator} are Python-centric, and the model itself is ecosystem-agnostic.

Node.js additionally hosts the orchestration command-line scripts under \texttt{scripts/ml}. These scripts coordinate the higher-level workflow steps, such as generating the repository seed, building the dataset, executing the training notebook, and staging artifacts. They wrap the underlying Python pipeline so that the full reproducible workflow can be triggered through consistent \texttt{npm} commands.

\section{Dataset Builder and Training Pipeline}
\label{appendix:impl-pipeline}

This section documents the offline path from repository seed to evaluated model
files: module layout, quality reporting, feature revision, model configuration,
artifact format, notebook orchestration, and the staging gate.

\subsection{Dataset-Builder Modules}
\label{sec:impl-dataset-modules}

This subsection expands the dataset builder of Section~\ref{sec:impl-dataset-builder}.

The historical dataset builder is divided into focused modules, summarized in Table~\ref{tab:impl-dataset-modules}, so that ingestion, event parsing, feature computation, and label computation can be tested independently.

\begin{table}[htbp]
\centering
\caption{Dataset builder modules}
\label{tab:impl-dataset-modules}
\begin{tabular}{p{0.22\textwidth} p{0.70\textwidth}}
\hline
\textbf{Module} & \textbf{Responsibility} \\
\hline
\texttt{adapters} & deps.dev, GitHub, and npm/PyPI registry adapters used for package-to-repository resolution and stable metadata. \\
\texttt{events} & GH Archive parsing into normalized repository histories of commits, issues, pull requests, releases, stars, and forks. \\
\texttt{features} & Observation-time feature computation from history before \(t\). \\
\texttt{labels} & 12-month maintenance and inactivity label computation from the future window. \\
\texttt{sampling} & Candidate sampling and popularity-tier derivation. \\
\texttt{pipeline} & The \texttt{DatasetBuilder} orchestrating ingestion, snapshots, labels, and export. \\
\texttt{storage} & JSON and JSON Lines reading and writing of the intermediate artifacts. \\
\texttt{cli} & The command-line entry point used by the orchestration scripts. \\
\hline
\end{tabular}
\end{table}

The \texttt{DatasetBuilder} class executes the construction workflow described in the methodology. The \texttt{ingest\_candidates} step loads and de-duplicates seed candidates, resolves each package to a GitHub repository, and writes the \texttt{repositories\sdot jsonl}, \texttt{packages.jsonl}, and \texttt{package\_repository\_links\sdot jsonl} files. Candidates that cannot be mapped to a GitHub repository, and forks when forks are excluded, are dropped at this stage. The \texttt{build\_snapshots} step iterates over quarterly observation dates and, for each repository, constructs an \texttt{ObservationSnapshot} with explicit feature and label windows before computing the observation-time feature row. The \texttt{build\_labels} step computes the future-window label for every snapshot, and \texttt{export\_training\_dataset} assembles the final training rows and the runtime repository feature cache.

\subsection{Dataset-Quality Reporting}
\label{sec:impl-dataset-quality}

To make the per-run documentation items required by the methodology (Section~\ref{sec:method-observation-window}) reproducible rather than transcribed by hand, a small \texttt{dataset\_quality} module computes a structured quality report from the exported snapshot rows. It produces three views. The first is row accounting: the total, labeled, and unlabeled counts, the inactive and active label balance, the number of already-archived-at-observation rows that are excluded from fitting, and the completeness signals introduced above, including the \texttt{repo\_silent\_before\_horizon} count that quantifies how many labeled examples the corrected coverage gate retains. The second is label balance per subgroup, broken down by ecosystem and by popularity tier, which exposes where labeled examples are concentrated. The third is per-feature missingness, derived from the missing-signal annotations already attached during feature extraction, so that a feature is reported as missing only when its underlying observation-time signal was unavailable rather than genuinely zero. The report is rendered in the training notebook alongside the existing coverage summary, and the row accounting deliberately records whether label-completeness signals are present in the export so that a value of zero is not misread as an absence of incomplete windows when an older export simply did not carry them.

\subsection{Feature-Set Refinement After the Formative Review}
\label{sec:impl-feature-refinement}

This subsection gives the details of the first design-cycle refinement summarized in Section~\ref{sec:impl-dataset-builder}.

The feature set used for training was refined iteratively, as described in the methodology's formative development step (Section~\ref{sec:method-iterative-development}). The first feature set contained variables that turned out to be poor point-in-time predictors. A representative example is a raw open-issue count: large, actively maintained projects can carry very large issue backlogs, so a model can wrongly associate a high absolute issue count with risk, even though the project is healthy. The revised feature set therefore replaces raw counts with relational and trend-based variables, such as issue closure ratio, issue backlog growth, stale-issue share, and pull-request response and merge latency, which are more comparable across projects of different sizes.

A second class of problems came from repository identity. Some packages point to mirror repositories rather than to the upstream development repository, which can make a project look inactive even though active development happens elsewhere. The mapping logic therefore prefers explicit seed and deps.dev mappings and resolves the primary package-to-repository link per repository, rather than treating every registry entry as an independent signal. Finally, the revised feature set keeps \texttt{repo\_archived\_at\_obs} only as diagnostic metadata and excludes already-archived observation rows during model fitting, so that the visible archival flag cannot become a shortcut label.

\subsection{Model Configurations and Hyperparameters}
\label{sec:impl-from-scratch}

This subsection expands the model families introduced in Section~\ref{sec:impl-models} and specified in Section~\ref{sec:method-model-training}.

The Logistic Regression model is implemented directly in Python. It standardizes each input feature using training-set means and standard deviations, stores these standardization parameters with the model, and fits the coefficients by gradient descent on the class-weighted log-likelihood. Class-balanced sample weights counter the label imbalance, and an L2 penalty regularizes the coefficients. Because the model stores explicit per-feature coefficients, it doubles as an interpretable baseline whose weights can be inspected to check whether the engineered features behave as expected. This interpretability is what made it useful during the feature-set refinement described in Appendix~\ref{sec:impl-feature-refinement}.

The neural network is a multilayer perceptron implemented on top of numpy, without a deep-learning framework. It uses two hidden layers of 64 and 32 units with ReLU activations and a sigmoid output, standardizes its inputs in the same way as the Logistic Regression model, and initializes its weights with Xavier-uniform values under a fixed random seed for reproducibility. Training uses mini-batch stochastic gradient descent on the binary cross-entropy objective, with balanced class weights and L2 regularization. The default configuration matches the methodology: a learning rate of \(0.01\), 150 epochs, a batch size of 256, and an L2 penalty of \(0.001\). The forward and backward passes are written explicitly, which keeps the model transparent and removes any heavy training dependency from the artifact.

The XGBoost model wraps the \texttt{xgboost} library. It trains a gradient-boosted tree ensemble with a binary-logistic objective and the same class-balanced sample weights as the other models. The default configuration uses 80 boosting rounds, a maximum tree depth of three, a learning rate of \(0.08\), row and column subsampling of \(0.9\), L2 leaf regularization, the histogram tree method, and a fixed seed. After training, the trained booster is serialized to its JSON representation and stored inside the model artifact, and the gain-based feature importances are computed and normalized so they can be displayed in the evaluation dashboard.

\subsection{Artifact Serialization}
\label{sec:impl-artifacts}

This subsection expands the artifact format referred to in Section~\ref{sec:impl-models}.

Each trained model is serialized to a self-describing JSON artifact. The artifact stores the model family and version, the ordered feature names, the model parameters (coefficients and standardization for Logistic Regression, weights and biases for the neural network, or the serialized booster for XGBoost), the calibration bins, the decision threshold, the feature-set version, and the training timestamp. Because the feature names and standardization are part of the artifact, the runtime service can construct the feature vector in the correct order and apply the same transformations without sharing code paths with the training process. Each training run is written to a run file under the training runs directory together with the dataset hash and the held-out metrics, and a \texttt{latest-run.json} pointer records the most recently produced run.

\subsection{Training Orchestration}
\label{sec:impl-training-orchestration}

This subsection expands the training workflow of Section~\ref{sec:impl-deployment}.

The headless command \texttt{npm run ml:train} executes the training notebook non-interactively and writes the run artifacts and the latest-run pointer, so that the same notebook is used for both interactive exploration and reproducible artifact generation. The end-to-end command \texttt{npm run ml:bootstrap} chains dataset construction and notebook-based training, and the foundation variants of these commands use the preserved 5{,}000-repository seed and the filtered GH Archive cache for the thesis-grade run. The orchestration scripts wrap the Python pipeline rather than reimplementing it; they pass through the seed file, the GH Archive source, the output directories, the observation window, and the dataset-size guardrails, which keeps the reproducibility parameters explicit at the command line.

\begin{lstlisting}[language=bash,caption={Headless notebook-based training},label={lst:impl-train}]
npm run ml:train
\end{lstlisting}

\subsection{Artifact Staging and Promotion Gate}
\label{sec:impl-staging}

This subsection expands the promotion gate summarized in Section~\ref{sec:impl-deployment}.

Trained artifacts are not deployed automatically. The staging command \texttt{npm run ml:stage\sdash training} validates a candidate training bundle and promotes it into the deployed bundle under \texttt{deployment\sslash training} only if it passes a set of quality gates. This gate is the safety boundary between experimentation and deployment.

The staging script first validates the dataset itself: the snapshots must be non-empty, every labeled snapshot must come from a real GitHub repository, and the bundle must contain at least a configured minimum number of distinct repositories and distinct inactive repositories. It then validates the model run artifacts: all required model artifacts must be present and complete for the latest dataset hash, and their feature versions must match the expected full-history and cold-start feature sets. This prevents a partial or inconsistent bundle, for example one missing the cold-start models, from reaching deployment.

The most important gate is the regression comparison. Before promotion, the candidate artifacts are compared against the currently deployed baseline. For each required model, the script checks that ROC-AUC has not dropped, and that the Brier score has not increased, by more than a configured tolerance. If any required model fails this comparison, promotion is rejected and the existing staged bundle is retained. When promotion succeeds, the script normalizes the snapshot file, the run artifacts, and the repository feature cache to their runtime paths and writes them into the deployed bundle.

\section{Runtime Services and Deployment}
\label{appendix:impl-runtime}

This section documents the components that consume the promoted model bundle at request time --- the scoring service, the backend that orchestrates an analysis, the analyst frontend, and the container and cluster deployment. None of them is part of the evidential argument of this thesis; no reported result depends on them.

\subsection{Scoring Service}
\label{sec:impl-scoring-service}

This subsection expands the scoring component of Section~\ref{sec:impl-architecture}, and implements the undefined-feature rules of Section~\ref{sec:method-undefined-features}.

The scoring service is a FastAPI application under \texttt{mltraining\sslash scoring\sslash app}. It exposes a small internal interface: health and readiness probes, a \texttt{\sslash features\sslash extract} endpoint for feature extraction, and the \texttt{/score/model} endpoint that scores a batch of dependencies with a supplied model artifact. The service does not own a model registry; instead, the backend passes the staged artifact in the scoring request, which keeps the scoring service stateless and makes the deployed model an explicit input rather than hidden state.

For each dependency, the service builds the feature vector in the artifact's feature order, imputing undefined signals from the training-set medians stored with the artifact (Section~\ref{sec:method-undefined-features}) and recording which signals were missing. It deserializes the artifact according to its model family, computes the raw probability, and applies the stored calibrator to obtain the calibrated inactivity probability. This probability is scaled to a 0--100 inactivity-risk score, and a complementary 12-month maintenance-outlook score is derived as its inverse. The service assigns a discrete risk bucket and an action level from the score, where the action level is one of \texttt{monitor}, \texttt{review}, or \texttt{replace\_candidate}, so that the continuous score maps onto the triage actions used in the demonstrator.

The service keeps the inactivity prediction and the security posture strictly separate, exactly as required by the methodology. The OpenSSF Scorecard signal is summarized into a parallel security-posture score and is never fed into the predictive feature vector \cite{zahan2023scorecard}.

Each score is accompanied by an \emph{evidence-support} value \(s \in [0,1]\). It is deliberately not a statistical confidence interval or a probability that the prediction is correct: the artifact ships no coefficient covariance, so no closed-form predictive variance is available. Instead, \(s\) is an operational measure of how much observable evidence backs the individual score. It is the geometric mean of three per-repository quantities:
\[
  s = \left( c \cdot d \cdot b \right)^{1/3},
\]
Here, \(c\) is signal coverage: the share of expected maintenance signals
observed before imputation, or the model's \texttt{signal\_completeness} feature
when present. The in-distribution share \(d\) is the fraction of observed
evidential features within \(2\sigma\) of the training mean. Finally,
\(b = n / (n + 30)\) maps the number of validation samples \(n\) backing the
matched calibration bin onto \([0,1]\). The geometric mean makes a single weak
component collapse the result. The decision margin to the threshold is reported
separately and is not part of \(s\). Explicit caveats are attached when
calibration bins are missing, when the artifact's feature version is unexpected,
or when input signals had to be imputed. Finally, the service returns
human-readable explanation factors and evidence items so that an analyst can see
which signals shaped the score rather than only the number itself.

\subsection{Backend Orchestration Service}
\label{sec:impl-backend}

This subsection expands the backend component of Section~\ref{sec:impl-architecture} and realizes the two prediction regimes of Section~\ref{sec:method-prediction-regimes}.

The backend is a Go service under \texttt{backend/api}. It exposes the public REST API for analyses, repositories, jobs, and the training dataset and run endpoints. Its central responsibility is to turn a submission into an enriched, scored analysis reliably, even in the presence of slow or failing external providers.

Analyses are processed through a durable job queue backed by PostgreSQL. When a submission arrives, the service creates an analysis record and a job record in a single transaction and returns immediately. A background worker leases the next pending job, processes it, and persists the result. If processing fails, the job is retried with a bounded number of attempts and a retry delay, and it is only marked as failed once the attempts are exhausted. This design makes the analysis lifecycle resilient and observable, and it decouples the user-facing request from the slower enrichment and scoring work. For repository submissions, the service also reuses a recent completed analysis when one exists for the same repository and the same staged model set, which avoids redundant re-enrichment.

During processing, the worker enriches each repository and scores it. Enrichment fetches stable repository metadata from GitHub and retrieves the OpenSSF Scorecard signal, attaching the raw signals to the repository for display. Each enrichment provider is implemented as a separate adapter, so a single failing provider degrades the available signals rather than the whole analysis.

The backend implements the two prediction regimes from the methodology at runtime. When it enriches a repository, it attempts to look the repository up in the staged repository feature cache. On a cache hit, the repository is scored in the full-history regime using the cached observation-time features; on a miss, it falls back to the cold-start regime using only approximate features derived from live metadata, and a caveat records this. When several model artifacts are staged for the same dataset hash and regime, the backend scores the repository with each and combines the results into an ensemble profile by averaging the scores, while also retaining the per-model outputs and adding a model-agreement explanation factor that reports the spread between models.

\subsection{Analyst Frontend}
\label{sec:impl-frontend}

This subsection expands the frontend component of Section~\ref{sec:impl-architecture}; the analyst workflow it serves is demonstrated in Chapter~\ref{chap:demonstrator}.

The frontend is a Next.js application under \texttt{frontend/web} that uses the React component model and the application router. It provides the analyst-facing workflow and the evaluation dashboards used during the thesis. Submission is handled through a form that accepts one or more repository URLs or a demo analysis. Once an analysis completes, the dashboard presents the repository overview with risk distribution and filtering, and a dedicated view renders the per-repository detail card with its risk profile, caveats, and evidence.

A second area of the frontend is the machine-learning evaluation surface. These pages read the training dataset and run endpoints exposed by the backend and visualize them with purpose-built charts, including a calibration curve, a model-metric comparison across the trained models, a Logistic Regression coefficient view, and a training-metric history. These views were used directly to inspect and report model behavior during the thesis.

\subsection{Container and Cluster Deployment}
\label{sec:appendix-impl-deployment}

This subsection expands the deployment path of Section~\ref{sec:impl-deployment}.

The artifact is containerized for both local and production-style operation. Canonical image definitions live under \texttt{deployment\sslash docker}, and a Docker Compose file orchestrates PostgreSQL, the Go backend, the scoring service, and the frontend for local development; a single \texttt{npm run dev} command builds and starts the full stack and waits for the services to become healthy. For production-style operation, the project is deployed to a k3s cluster managed by Argo CD, with container images published to a registry and the application served behind a public domain \cite{ossRiskRadarReadme}. Because the deployed model bundle is a versioned set of files under \texttt{deployment\sslash training}, deployment and model promotion remain auditable: the running service uses exactly the artifacts that passed the staging gate described in Appendix~\ref{sec:impl-staging}.

\section{Evaluation Instruments}
\label{appendix:impl-eval-instruments}

This section gives the mechanics of the two evaluation instruments introduced in Section~\ref{sec:impl-models}. They are of different kinds, and the distinction matters for what each result may claim. The slice evaluation is a re-analysis: it re-scores the predictions of the same fitted model within subgroups. The repository-disjoint split is a separate training run under a different partitioning rule, so it produces its own fitted model. Neither alters the reported run, whose serialized artifacts and headline metrics are unchanged, because the reported configuration retains the chronological split.

\subsection{Subgroup Definitions for the Slice Evaluation}
\label{sec:impl-slice-evaluation}

The resulting figures are reported in Section~\ref{sec:results-slices} and tabulated in Appendix~\ref{appendix:eval-tables}.

A single aggregate metric can hide a subgroup on which the model performs poorly, because a strong majority subgroup can mask a weak minority one. A slice-evaluation helper therefore re-scores the same held-out predictions within subgroups, reusing the existing metric implementation so that no separate evaluation logic is introduced. To make this possible the training pipeline also returns the per-row held-out predictions alongside the aggregate result; these are used only for analysis and are not written into the model artifact, so the serialized artifacts and the reported headline metrics are unchanged. The harness includes a consistency check that requires the aggregate over all slices to reproduce the model's reported held-out ROC-AUC, so the per-slice figures are guaranteed to originate from the same execution as the headline metrics.

The subgroups are those that vary meaningfully in the foundation dataset: popularity tier, a cold-start-like versus has-history regime, and whether the repository was still active at the observation date. The last of these is the active-at-\(t\) subset referenced in the evaluation strategy (Section~\ref{sec:method-evaluation-strategy}) and in the discussion of activity persistence in Chapter~\ref{chap:validity}: restricting the evaluation to repositories that are still active at \(t\) separates genuine early-warning skill from the trivial tendency of already-quiet repositories to remain quiet. Two candidate dimensions are deliberately excluded because the sampling frame holds them constant. Ecosystem is excluded because the repository-first foundation seed assigns the same provenance ecosystem to every row; cross-ecosystem evaluation would require a package-first sampling frame and is left to future work. Repository-age band is excluded because the offline build carries no repository creation date, so \texttt{repo\_age\_days} is identically zero and every row would fall into a single band (Section~\ref{sec:method-offline-metadata}); reporting such a slice would suggest an age analysis that the data cannot support. Because slicing reduces the sample size, each slice is reported with its own row count, ranking metrics that are undefined for single-class slices are reported as absent rather than as a misleading value, and thin slices are interpreted with corresponding caution.

\subsection{Mechanics of the Repository-Disjoint Split}
\label{sec:impl-grouped-split}

This subsection addresses the repository-recurrence threat of Chapter~\ref{chap:validity} that arises from the split described in Section~\ref{sec:method-observation-window}.

The reported evaluation uses a chronological split: rows are ordered by observation date and cut into training, validation, and test partitions. Because observations are taken quarterly, a repository contributes several snapshots over time, so the same repository can appear in both the training and the test partition, and the twelve-month label windows of its successive snapshots overlap. The held-out metric then partly reflects generalization to later observations of repositories already seen during training, which is optimistic relative to performance on genuinely new repositories. This is the deployment-realistic question --- the system does re-score repositories over time --- but it should not be the only question asked.

A further evaluation iteration therefore adds a repository-disjoint split that assigns every snapshot of a repository to a single partition, so that the held-out set contains only repositories absent from training. The split is selectable through a configuration option whose default preserves the chronological behavior, so the reported artifacts and headline metrics are unchanged, and the pipeline asserts that the partitions are repository-disjoint whenever the grouped split is used.

Because grouping changes which rows are available for fitting, this is a full second training run rather than a re-scoring of the reported model. The same pipeline entry point is invoked with the grouped strategy, and it refits in the same order as the reported run: the feature preprocessing, then the classifier on the grouped training partition, then the histogram calibrator and the F1-selected decision threshold on the grouped validation partition. The grouped test partition enters only after those decisions are fixed. The two runs therefore report two different fitted models, and the difference between their held-out ROC-AUC values bounds the optimism introduced by repository recurrence. In practice the held-out ROC-AUC decreases only marginally, which indicates that the model's discrimination does not rest on memorizing individual repositories. The comparison is between held-out repositories drawn from the same retrospectively sampled foundation frame (Section~\ref{sec:method-data-sources}), so it does not establish transfer to open-source repositories in general. The exact figures for both splits are produced in the training notebook and reported in Chapter~\ref{chap:evaluation}.

\section{Dataset Construction Steps}
\label{appendix:dataset-construction}

This section expands the dataset construction workflow of Section~\ref{sec:method-dataset-construction}; its software realization is described in Section~\ref{sec:impl-dataset-builder}.

The historical dataset builder creates its dataset through the sequence of intermediate steps summarized in Table~\ref{tab:method-dataset-flow}, writing the intermediate files listed below under the configured output directory so that the construction process is auditable and the generated dataset can be inspected before training.

\begin{table}[htbp]
\centering
\caption{Historical dataset construction workflow}
\label{tab:method-dataset-flow}
\begin{tabular}{p{0.08\textwidth} p{0.82\textwidth}}
\hline
\textbf{Step} & \textbf{Description} \\
\hline
1 & Load repository candidates from the foundation seed. \\
2 & Use the seed buckets to obtain a mixture of active, dormant, and archived repository histories. \\
3 & Resolve repository identities from explicit seed fields or package metadata where required. \\
4 & Enrich stable repository metadata through GitHub. \\
5 & Parse GH Archive event files into normalized repository histories. \\
6 & Build quarterly observation snapshots for each repository. \\
7 & Compute observation-time features using only data available at or before \(t\). \\
8 & Compute 12-month labels using only the future window \((t, t + 12\text{ months}]\). \\
9 & Export the final snapshot dataset for the existing model-training workflow. \\
10 & Write a repository feature cache for runtime full-history scoring. \\
\hline
\end{tabular}
\end{table}

% Nearly every word here is monospace, whose interword glue cannot stretch, so TeX has
% nothing to give and would otherwise pick an overfull line. sloppypar permits loose ones.
\begin{sloppypar}
The documented intermediate outputs are \texttt{repositories\sdot jsonl}, \texttt{packages\sdot jsonl}, \texttt{package\su repository\su links\sdot jsonl}, \texttt{observation\su snapshots\sdot jsonl}, \texttt{snapshot\su features\sdot jsonl}, \texttt{snapshot\su labels\sdot jsonl}, \texttt{training\su snapshots\sdot json}, and \texttt{repository\sdash feature\sdash cache\sdot json}.
\end{sloppypar}

\section{Evaluation Metric Definitions}
\label{sec:method-metric-definitions}

This section gives the formal definitions of the metrics named in the evaluation
strategy (Section~\ref{sec:method-evaluation-strategy}) and reported throughout
Chapter~\ref{chap:evaluation}.

All metrics are computed on the calibrated probabilities of the held-out test
set. Let \(N\) be the number of evaluation samples, \(p_i \in [0,1]\) the
predicted probability of inactivity for sample \(i\), and \(y_i \in \{0,1\}\) the
true label (\(1\) = inactive within twelve months). The positive (inactive) base
rate of the test set is
\begin{equation}
  \bar{p} = \frac{1}{N}\sum_{i=1}^{N} y_i ,
\end{equation}
reported as the \emph{inactive 12m rate}.

\paragraph{Ranking.} ROC-AUC (equivalently \(\mathrm{AUROC}\), the notation used
in the equation below) is computed as the Mann--Whitney statistic over all
positive--negative pairs, that is, the probability that a random inactive
repository is scored above a random active one, with ties counted as one half:
\begin{equation}
  \mathrm{AUROC} = \frac{1}{N_{+}N_{-}}
    \sum_{i:\,y_i=1}\;\sum_{j:\,y_j=0}
    \Big( \mathbb{1}[p_i > p_j] + \tfrac{1}{2}\,\mathbb{1}[p_i = p_j] \Big),
\end{equation}
where \(N_{+}\) and \(N_{-}\) are the numbers of positive and negative samples.
It is threshold- and calibration-independent: \(0.5\) is random ranking and
\(1.0\) is a perfect ranking.

\paragraph{Ranking under class imbalance.} Average precision summarizes the
precision-recall curve. Samples are sorted by decreasing score; at each distinct
score threshold \(n\), let \(P_n\) and \(R_n\) be the precision and recall
attained by predicting positive down to that threshold. Average precision is the
step-wise sum
\begin{equation}
  \mathrm{AP} = \sum_{n} \big( R_n - R_{n-1} \big)\, P_n , \qquad R_0 = 0 ,
\end{equation}
with tied scores collapsed into a single threshold so that a signal with many
ties is neither rewarded nor penalized for them. No interpolation between points
is applied. Unlike ROC-AUC, average precision ignores true negatives and so does
not credit a model for correctly ranking a large negative class. Its no-skill
reference is not \(0.5\) but the positive base rate \(\bar{p}\), which a random
ranking attains, so an AP value is interpretable only alongside the prevalence of
the set it was computed on. It is reported for the early-warning subset in
Section~\ref{sec:results-precision-recall}.

\paragraph{Probabilistic accuracy.} The Brier score is the mean squared error
between probability and outcome, and the log loss is the binary cross-entropy
(with a small \(\varepsilon\) for numerical stability):
\begin{align}
  \mathrm{Brier}   &= \frac{1}{N}\sum_{i=1}^{N} (p_i - y_i)^2, \\
  \mathrm{LogLoss} &= -\frac{1}{N}\sum_{i=1}^{N}
     \Big[\, y_i \ln p_i + (1-y_i)\ln(1-p_i) \,\Big].
\end{align}
Both are lower-is-better and penalize miscalibration; the log loss grows without
bound as a confident prediction approaches the wrong outcome, so it punishes
overconfidence more sharply than the bounded Brier score.

\paragraph{Calibration.} The expected calibration error partitions
\([0,1]\) into \(B = 10\) equal-width bins. For bin \(b\) with member set
\(\mathcal{B}_b\), let \(\mathrm{conf}_b\) be the mean predicted probability and
\(\mathrm{acc}_b\) the empirical positive rate; the ECE is the sample-weighted
average gap between predicted confidence and observed frequency:
\begin{equation}
  \mathrm{ECE} = \sum_{b=1}^{B} \frac{|\mathcal{B}_b|}{N}\,
    \big|\,\mathrm{conf}_b - \mathrm{acc}_b\,\big|.
\end{equation}

\paragraph{Thresholded classification.} A decision threshold \(\tau\) maps
probabilities to decisions \(\hat{d}_i = \mathbb{1}[p_i \ge \tau]\), yielding the
counts of true and false positives and negatives (\(\mathrm{TP}, \mathrm{FP},
\mathrm{FN}, \mathrm{TN}\)). From these,
\begin{equation}
  \mathrm{Precision} = \frac{\mathrm{TP}}{\mathrm{TP}+\mathrm{FP}}, \qquad
  \mathrm{Recall} = \frac{\mathrm{TP}}{\mathrm{TP}+\mathrm{FN}}, \qquad
  \mathrm{F1} = \frac{2\,\mathrm{Precision}\cdot\mathrm{Recall}}
                    {\mathrm{Precision}+\mathrm{Recall}} .
\end{equation}
The F1 score is the harmonic mean of precision and recall, so it cannot be made
large by optimizing one at the expense of the other; \(\tau\) is selected on the
validation set to maximize it.

\paragraph{Combined quality score.} The quality score converts ranking and
probabilistic accuracy into skill scores relative to the trivial base-rate
predictor and averages them, weighting ranking at \(0.6\) and calibration at
\(0.4\). With the climatological Brier score \(\bar{p}(1-\bar{p})\) and
\(\operatorname{clip}_{[0,1]}\) clamping to the unit interval,
\begin{align}
  s_{\mathrm{AUC}}   &= \operatorname{clip}_{[0,1]}\!\left(\frac{\mathrm{AUROC} - 0.5}{0.5}\right), \\
  s_{\mathrm{Brier}} &= \operatorname{clip}_{[0,1]}\!\left(1 - \frac{\mathrm{Brier}}{\bar{p}(1-\bar{p})}\right), \\
  \mathrm{Quality}   &= \operatorname{clip}_{[0,1]}\!\big(0.6\,s_{\mathrm{AUC}} + 0.4\,s_{\mathrm{Brier}}\big).
\end{align}
A value of \(0\) indicates no improvement over predicting the base rate and \(1\)
indicates perfect ranking and calibration. Because both skill terms are measured
against the test-set base rate \(\bar{p}\), the quality score is conditional on the
sampled prevalence and is intended for model comparison rather than as a
standalone performance guarantee.

\section{Reproducibility}
\label{sec:method-reproducibility}

This section makes the dataset construction workflow of Section~\ref{sec:method-dataset-construction} reproducible, and records the identifying metadata of the run whose results are reported in Chapter~\ref{chap:evaluation}.

The dataset builder and training pipeline are executed through repository scripts. The foundation seed can be generated with:

\begin{verbatim}
npm run ml:seed:foundation -- `
  --target-repositories 5000 `
  --github-token <TOKEN>
\end{verbatim}

The larger foundation dataset can then be built and trained with:

\begin{verbatim}
npm run ml:bootstrap:foundation -- `
  --github-token <TOKEN> `
  --gharchive-source <GHARCHIVE_PATH>
\end{verbatim}

The underlying full command path uses a generated foundation seed, a GH Archive source directory, a dedicated training output directory, a training snapshot export path, and a repository feature cache output path. The repository documentation recommends the foundation path for the thesis/training base, while the smaller default bootstrap path remains a local smoke-test workflow \cite{ossRiskRadarDatasetBuilder}.

For each final training run, the following items should be documented:

\begin{itemize}
    \item seed generation command and seed metadata,
    \item GH Archive coverage period,
    \item output directory,
    \item training snapshot path,
    \item repository feature cache path,
    \item dataset hash,
    \item model artifact path,
    \item model name and version,
    \item feature version,
    \item number of total, labeled, and unlabeled rows,
    \item train, validation, and test split sizes,
    \item held-out evaluation metrics.
\end{itemize}

The reported full-history run uses dataset hash:

\begin{verbatim}
2d29a928dac34ea1dab919f87e89386f2e6eac3bda59aee7a625c5ad0873b0d8
\end{verbatim}

and exports the model artifact:

\begin{verbatim}
deployment/training/runs/
  20260709T205727480957Z-xgboost-full-history-2d29a928dac3.json
\end{verbatim}

The reported held-out metrics are those recorded in this staged run file, which is the authoritative record of the reported experiment. The corresponding source revision --- containing this dataset, the staged run artifacts, and the training notebook --- is archived at commit \texttt{f5e688e} \cite{ossRiskRadarReadme}, so that the reported run can be traced to a fixed state of the dataset builder and training pipeline. One caveat applies to the hash above: it is the value recorded by the staged run at training time, whereas the dataset file committed at \texttt{f5e688e} was re-serialized afterwards and therefore carries a different content hash (\texttt{f71b9f7a\ldots}). Re-training \texttt{xgboost\sdash full\sdash history} on the committed file reproduces the recorded held-out metrics exactly --- ROC-AUC \(0.957852\), Brier \(0.077939\), expected calibration error \(0.035155\), F1 \(0.916369\), and the selected threshold \(0.506689\), each identical to six decimal places --- so the two files differ only in serialization: the recorded hash identifies the run, and the committed file reproduces it.\kiref{Claude Opus 5/9}

\section{Label-Definition Sensitivity}
\label{appendix:label-sensitivity}

Table~\ref{tab:label-sensitivity} reports the robustness check of
Section~\ref{sec:results-label-sensitivity}: the held-out evaluation repeated under
one-of-four, two-of-four (primary), and three-of-four variants of the maintenance
rule. All \(50{,}000\) snapshots are re-labeled from the same stored future-window
signals, and the full-history XGBoost model is re-trained per variant on the
identical time-aware split. Because the recency baseline is undefined for the
\(444\) test rows without an observable commit history
(Section~\ref{sec:results-baselines}), the model is reported both on the full test
set and re-scored on the \(4{,}556\) recency-defined rows, so that the model and
the two trivial rules are compared on identical rows. Two checks confirm that the
re-labeling harness matches the primary pipeline: the two-of-four labels derived
from the stored signals agree with the exported \texttt{label\_inactive\_12m} on
all \(50{,}000\) snapshots, and the two-of-four row reproduces the reported
headline figures. The analysis is reproduced by
\texttt{scripts\sslash ml\sslash label\su sensitivity\sdot py}, which performs
both checks and aborts if the first one fails
(Appendix~\ref{appendix:code-availability}).

\clearpage
\begin{table}[htbp]
  \centering
  \small
  \setlength{\tabcolsep}{3.5pt}
  \caption{Sensitivity of the held-out evaluation to the maintenance-rule
  threshold. ``Inactive rate'' is the positive-class share of the test set.
  ``XGBoost (all)'' is the full-history model on all \(5{,}000\) test rows; the
  remaining columns are on the \(4{,}556\) recency-defined rows, where model and
  baselines are scored on identical rows. ``Margin'' is the model there minus
  the stronger of the two trivial rules, which is the recency rule under the
  permissive variant and commit volume under the other two.}
  \label{tab:label-sensitivity}
  \begin{tabular}{lrrrrrr}
    \toprule
    Maintenance rule & Inactive rate & XGBoost (all) & XGBoost (rec.) & Recency & Commits & Margin \\
    \midrule
    1-of-4 (permissive) & 0.581 & 0.961 & 0.975 & 0.965 & 0.960 & \(+0.010\) \\
    2-of-4 (primary)    & 0.661 & 0.958 & 0.974 & 0.957 & 0.958 & \(+0.016\) \\
    3-of-4 (strict)     & 0.746 & 0.959 & 0.972 & 0.940 & 0.949 & \(+0.023\) \\
    \bottomrule
  \end{tabular}
\end{table}
\section{Full Demonstrator Ranking}
\label{appendix:demo-full}

Table~\ref{tab:demo-ranking-full} lists the complete ranked inactivity-risk view for all 30 dependency repositories of \texttt{ArchipelagoMW\sslash Archipelago} produced by the model in the demonstrator (Chapter~\ref{chap:demonstrator}); the body reports the flagged and medium-risk subset in Table~\ref{tab:demo-ranking}.

The explanation factors of the highest-ranked repository, \texttt{bsdiff4} (\texttt{ilanschnell\sslash bsdiff4}, risk \(92.1\); Section~\ref{sec:demo-ranking}), taken directly from the scored output, illustrate the intended semantics: the two cold-start ensemble members predict \(97.0\%\) and \(87.2\%\) inactivity risk, the model-agreement factor reports the spread between them, and the input-completeness factor records that \(75\%\) of expected signals were available before imputation. The moderate evidence support of \(0.69\) and the caveats attached to the score --- that it came from the cold-start model, no full-history cache row being available, and that two release-related signals were imputed --- are what present the flag as a prompt for manual review rather than a definitive judgement.

\begin{table}[htbp]
  \centering
  \small
  \caption{Complete ranked inactivity-risk view for the dependency repositories of
  \texttt{ArchipelagoMW/Archipelago}, as produced by the model. Risk is
  the 0--100 inactivity-risk score; ``Supp.'' is the per-prediction evidence support;
  ``Sec.'' is the parallel 0--100 security-posture score derived from the
  repository's OpenSSF Scorecard record, which is not an input to the inactivity
  model (Section~\ref{sec:method-research-design}) and is shown as --- for the three
  repositories with no available Scorecard record; ``Regime'' is full-history (FH)
  or cold-start (CS).}
  \label{tab:demo-ranking-full}
  \footnotesize
  \setlength{\tabcolsep}{3.5pt}
  \begin{tabular}{llrlrrl}
    \toprule
    Package & Source repository & Risk & Action & Supp. & Sec. & Regime \\
    \midrule
    bsdiff4 & \texttt{ilanschnell/bsdiff4} & 92.1 & \texttt{replace\_candidate} & 0.69 & 6 & CS \\
    setproctitle & \texttt{dvarrazzo/py-setproctitle} & 80.4 & \texttt{replace\_candidate} & 0.61 & 11 & CS \\
    \midrule
    jinja2 & \texttt{pallets/jinja} & 77.1 & review & 0.77 & 58 & CS \\
    cymem & \texttt{explosion/cymem} & 74.9 & review & 0.75 & 24 & CS \\
    markupsafe & \texttt{pallets/markupsafe} & 74.2 & review & 0.75 & 54 & CS \\
    pymem & \texttt{srounet/Pymem} & 54.3 & review & 0.47 & --- & CS \\
    \midrule
    flask-compress & \texttt{colour-science/flask-compress} & 48.5 & monitor & 0.45 & 13 & CS \\
    pony & \texttt{ponyorm/pony} & 42.8 & monitor & 0.59 & 1 & CS \\
    flask-limiter & \texttt{alisaifee/flask-limiter} & 41.9 & monitor & 0.58 & 18 & CS \\
    kivymd & \texttt{kivymd/KivyMD} & 41.8 & monitor & 0.47 & 15 & CS \\
    waitress & \texttt{Pylons/waitress} & 41.3 & monitor & 0.59 & 30 & CS \\
    schema & \texttt{keleshev/schema} & 37.3 & monitor & 0.78 & 23 & CS \\
    pyyaml & \texttt{yaml/pyyaml} & 35.4 & monitor & 0.77 & 43 & CS \\
    orjson & \texttt{ijl/orjson} & 31.1 & monitor & 0.61 & 27 & CS \\
    colorama & \texttt{tartley/colorama} & 27.8 & monitor & 0.54 & 30 & CS \\
    pyshortcuts & \texttt{newville/pyshortcuts} & 25.0 & monitor & 0.75 & 11 & CS \\
    werkzeug & \texttt{pallets/werkzeug} & 24.3 & monitor & 0.76 & 58 & CS \\
    certifi & \texttt{certifi/python-certifi} & 23.5 & monitor & 0.57 & 59 & CS \\
    platformdirs & \texttt{tox-dev/platformdirs} & 22.6 & monitor & 0.77 & --- & CS \\
    pathspec & \texttt{cpburnz/python-pathspec} & 16.6 & monitor & 0.83 & --- & CS \\
    kivy & \texttt{kivy/kivy} & 15.0 & monitor & 0.85 & 36 & CS \\
    websockets & \texttt{python-websockets/websockets} & 10.3 & monitor & 0.90 & 52 & CS \\
    flask & \texttt{pallets/flask} & 9.5 & monitor & 0.91 & 77 & FH \\
    flask-cors & \texttt{corydolphin/flask-cors} & 8.9 & monitor & 0.91 & 46 & CS \\
    typing-extensions & \texttt{python/typing\_extensions} & 5.8 & monitor & 0.95 & 99 & CS \\
    bokeh & \texttt{bokeh/bokeh} & 5.3 & monitor & 0.71 & 65 & CS \\
    flask-caching & \texttt{pallets-eco/flask-caching} & 5.0 & monitor & 0.95 & 49 & CS \\
    mistune & \texttt{lepture/mistune} & 4.2 & monitor & 0.96 & 43 & CS \\
    psutil & \texttt{giampaolo/psutil} & 3.4 & monitor & 0.72 & 47 & CS \\
    cython & \texttt{cython/cython} & 2.6 & monitor & 0.98 & 31 & CS \\
    \bottomrule
  \end{tabular}
\end{table}

\clearpage
\section{Per-Subgroup Metrics of the Slice Evaluation}
\label{appendix:eval-tables}

This section gives the full per-subgroup breakdown of \texttt{xgboost\sdash full\sdash history} whose three findings are interpreted in Section~\ref{sec:results-slices}.

\begin{table}[htbp]
  \centering
  \caption{Held-out ROC-AUC and Brier score of \texttt{xgboost-full-history}
  within subgroups of the test set. ``Inactive rate'' is the positive rate in
  the subgroup.}
  \label{tab:slice-metrics}
  \begin{tabular}{llrrr}
    \toprule
    Dimension & Subgroup & \(n\) & Inactive rate & ROC-AUC \\
    \midrule
    Popularity      & high            & 2297 & 0.527 & 0.972 \\
    Popularity      & medium          & 1753 & 0.778 & 0.964 \\
    Popularity      & low             & 950  & 0.767 & 0.846 \\
    \midrule
    Activity at \(t\) & active@\(t\)   & 1794 & 0.225 & 0.905 \\
    Activity at \(t\) & quiet@\(t\)    & 3206 & 0.905 & 0.911 \\
    \midrule
    History regime  & has-history     & 2295 & 0.350 & 0.921 \\
    History regime  & cold-start-like & 2705 & 0.924 & 0.929 \\
    \bottomrule
  \end{tabular}
\end{table}

\section{Ablation: The Human-Actor Filter}
\label{appendix:bot-ablation}

Commit- and contributor-based features and the inactivity label are computed from
human actors only; automated (bot) accounts are excluded
(Section~\ref{sec:method-feature-engineering}). Because this design choice could
plausibly be omitted, it was tested by rebuilding the full 5{,}000-repository
dataset with the filter disabled through a single build flag, with every other
setting identical, yielding 50{,}000 snapshots that align one-to-one with the
filtered build. The headline outcome is summarized in
Section~\ref{sec:results-label-sensitivity}: the filter earns its place on
construct-validity grounds rather than on discrimination. Three effects were
measured.

\paragraph{Label corruption.} Disabling the filter flips 357 of 50{,}000 labelled snapshots (\(0.71\%\)) from inactive to active, and \emph{none} in the opposite direction. The effect is strictly one-directional: a bot commit or merge inside the future window can satisfy the two-of-four maintenance rule, so an otherwise dormant repository is relabelled as maintained, whereas bots never make an active repository look inactive. Omitting the filter would therefore bias the target by systematically \emph{masking} inactivity, which is the error direction that matters most for a risk signal. That asymmetry is invisible to a headline metric computed on the corrupted data, which is why the filter is retained despite the small predictive effect below.

\paragraph{Feature inflation.} On the observation-time side, bots inflate the activity features substantially: with the filter disabled, \texttt{commits\_365d} changes on 7{,}858 of 50{,}000 snapshots (\(15.7\%\)) by a mean of \(+938\) commits, though the median change is only \(+80\), so the mean is driven by a small number of heavily automated repositories. \texttt{commits\_90d} rises by a mean of \(+388\) and \texttt{contributors\_365d} by \(+1.4\) distinct contributors, while pull-request \emph{count} features are unaffected because they are not actor-filtered. A small number of bot-driven repositories would thus dominate the activity signal.

\paragraph{Predictive effect.} A naive comparison, in which each model is scored against the labels of the dataset it was trained on, would show bots to be nearly harmless --- but it is misleading, because the unfiltered model is then scored against a partially corrupted target. A cleaner test trains on the bot-contaminated \emph{features} while evaluating against the clean, human-filtered \emph{labels} (Table~\ref{tab:bot-ablation-crosstest}): the contaminated features are consistently worse across all four families, though by a small margin (ROC-AUC falls by \(0.0006\)--\(0.0011\); Brier rises by up to \(0.0010\)).

\begin{table}[htbp]
  \centering
  \caption{Cross-test of the human-actor filter: training on bot-contaminated
  features versus clean features, both evaluated against the identical clean
  (human-filtered) label on the test set. A lower ROC-AUC or higher Brier
  for the contaminated features means the filter helps.}
  \label{tab:bot-ablation-crosstest}
  \small
  \setlength{\tabcolsep}{5pt}
  \begin{tabular}{lrrrr}
    \toprule
    & \multicolumn{2}{c}{ROC-AUC} & & \\
    \cmidrule(lr){2-3}
    Model & clean & contam. & \(\Delta\)ROC-AUC & \(\Delta\)Brier \\
    \midrule
    \texttt{logistic-regression-full-history} & 0.9494 & 0.9484 & \(-0.0010\) & \(+0.0005\) \\
    \texttt{xgboost-full-history}             & 0.9579 & 0.9569 & \(-0.0010\) & \(+0.0001\) \\
    \texttt{logistic-regression-cold-start}   & 0.9318 & 0.9308 & \(-0.0011\) & \(+0.0005\) \\
    \texttt{xgboost-cold-start}               & 0.9531 & 0.9524 & \(-0.0006\) & \(+0.0010\) \\
    \bottomrule
  \end{tabular}
\end{table}
